\documentclass[11pt]{article}
\usepackage[margin=1in]{geometry}
\usepackage{amsmath,amssymb,bm}
\usepackage{physics}
\usepackage{hyperref}

\title{Classical Single-Mode Cavity--Molecule Hamiltonians:\\
Direct Laser Driving vs Cavity-Mediated Driving}
\author{}
\date{}

\begin{document}
\maketitle

\section{General setup}

We consider a molecular system with classical nuclear degrees of freedom
\[
\bm R=\{\bm r_i\}_{i=1}^N, \qquad
\bm P=\{\bm p_i\}_{i=1}^N,
\]
with masses $m_i$ and an internal molecular potential $U(\bm R)$.
The system is coupled to a \emph{single electromagnetic cavity mode}, represented
classically as a harmonic oscillator with coordinate $q$ and momentum $p$:
\begin{equation}
H_{\text{cav}}(q,p)=\frac{p^2}{2}+\frac{\omega_c^2 q^2}{2},
\end{equation}
where $\omega_c$ is the cavity resonance frequency.

The molecular dipole projected along the cavity polarization direction
is denoted by the scalar function $\mu(\bm R)$.

\paragraph{Light--matter coupling (redefined convention).}
We adopt the linear interaction
\begin{equation}
H_{\text{int}}(\bm R,q)= -\lambda\,\omega_c\, q\,\mu(\bm R),
\label{eq:Hint}
\end{equation}
where $\lambda$ is a dimensionless coupling constant under the convention
that $q$ is proportional to the intracavity electric-field amplitude.

\paragraph{Optional dipole self-energy.}
A quadratic term consistent with the above scaling is
\begin{equation}
H_{\text{self}}(\bm R)=\frac{\lambda^2}{2}\,\mu(\bm R)^2.
\label{eq:Hself}
\end{equation}
This term stabilizes the energy landscape at strong coupling and is often
retained in length-gauge formulations; it may be omitted depending on the
approximation used.

\paragraph{Baseline Hamiltonian (no external driving).}
\begin{equation}
H_0=
\sum_{i=1}^N \frac{\bm p_i^2}{2m_i}
+U(\bm R)
+\frac{p^2}{2}+\frac{\omega_c^2 q^2}{2}
-\lambda\,\omega_c\, q\,\mu(\bm R)
+\frac{\lambda^2}{2}\mu(\bm R)^2.
\label{eq:H0}
\end{equation}

\section{Case (1): External laser drives the molecules directly}

In this scenario, the cavity mode is present but the external laser field
couples \emph{directly} to the molecular dipole. The laser field is prescribed
and does not respond to the molecular or cavity dynamics:
\begin{equation}
E_{\text{ext}}(t)=E_0\,s(t)\cos(\omega_L t+\phi),
\end{equation}
where $\omega_L$ is the laser frequency and $s(t)$ is a smooth envelope.

The laser--molecule interaction is
\begin{equation}
H_{\text{mol--laser}}(t)= -\mu(\bm R)\,E_{\text{ext}}(t).
\end{equation}

The full Hamiltonian is
\begin{equation}
H^{(1)}(t)=H_0-\mu(\bm R)\,E_{\text{ext}}(t).
\end{equation}

\subsection*{Equations of motion (case 1)}

\paragraph{Nuclear motion.}
\begin{equation}
\dot{\bm r}_i=\frac{\bm p_i}{m_i},
\end{equation}
\begin{align}
\dot{\bm p}_i &=
-\nabla_{\bm r_i}U(\bm R)
+\lambda\,\omega_c\, q\,\nabla_{\bm r_i}\mu(\bm R)
-\lambda^2\,\mu(\bm R)\,\nabla_{\bm r_i}\mu(\bm R)
\nonumber\\
&\quad
+E_{\text{ext}}(t)\,\nabla_{\bm r_i}\mu(\bm R),
\end{align}
where the term proportional to $\lambda^2$ is omitted if $H_{\text{self}}$
is not included.

\paragraph{Cavity mode.}
\begin{align}
\dot q &= p,\\
\dot p &= -\omega_c^2 q + \lambda\,\omega_c\,\mu(\bm R).
\end{align}

\paragraph{Physical note.}
In this case, the laser field directly forces the molecular dipole at frequency
$\omega_L$. The cavity does \emph{not} filter the laser spectrum: any frequency
component present in $E_{\text{ext}}(t)$ acts directly on the nuclei.
The cavity responds only indirectly through its coupling to $\mu(\bm R)$.

\section{Case (2): External laser pumps the cavity mode only}

Here the external laser is treated as an \emph{input field} that injects energy
into the cavity through its boundaries (mirrors or waveguide).
The laser does \emph{not} directly couple to the molecular dipole.

The cavity drive is represented by a time-dependent generalized force:
\begin{equation}
H_{\text{cav--drive}}(t)= - f(t)\,q,
\end{equation}
with
\begin{equation}
f(t)=f_0\,s(t)\cos(\omega_d t+\phi),
\end{equation}
where $\omega_d$ is the external driving frequency.

The full Hamiltonian is
\begin{equation}
H^{(2)}(t)=H_0 - f(t)\,q.
\end{equation}

\subsection*{Equations of motion (case 2)}

\paragraph{Nuclear motion.}
\begin{equation}
\dot{\bm r}_i=\frac{\bm p_i}{m_i},
\end{equation}
\begin{equation}
\dot{\bm p}_i=
-\nabla_{\bm r_i}U(\bm R)
+\lambda\,\omega_c\, q\,\nabla_{\bm r_i}\mu(\bm R)
-\lambda^2\,\mu(\bm R)\,\nabla_{\bm r_i}\mu(\bm R),
\end{equation}
again omitting the last term if $H_{\text{self}}$ is not included.

\paragraph{Cavity mode.}
\begin{align}
\dot q &= p,\\
\dot p &= -\omega_c^2 q + \lambda\,\omega_c\,\mu(\bm R) + f(t).
\end{align}
Equivalently,
\begin{equation}
\ddot q + \omega_c^2 q = \lambda\,\omega_c\,\mu(\bm R) + f(t).
\end{equation}

\paragraph{Physical note (filtered driving).}
In this case, the cavity acts as a \emph{frequency filter}.
Although the external drive oscillates at frequency $\omega_d$,
the intracavity field $q(t)$ responds most strongly when
$\omega_d \approx \omega_c$.
Off-resonant frequency components are suppressed, and the molecule
experiences a field shaped by the cavity response rather than the
raw laser spectrum.
This filtering, together with the dynamical back-action
between $q(t)$ and $\mu(\bm R)$, is the defining feature of cavity-mediated
driving.

\section{Conceptual comparison}

\begin{center}
\begin{tabular}{lcc}
\hline
 & Direct molecular drive & Cavity-mediated drive \\
\hline
Laser frequency seen by molecule & $\omega_L$ (unfiltered) & $\omega_c$-filtered response \\
Field degree of freedom & Prescribed & Dynamical ($q,p$) \\
Back-action from molecule & None & Present \\
Energy storage / memory & None & Yes (in cavity mode) \\
\hline
\end{tabular}
\end{center}

\end{document}
